\begin{Solbook}{1}
 \par \noindent  \begin {listEX}[3] \item $-256 + 32 - 744 \\= (-256 -744) + 32 \\= -1000 + 32 = -968$; \item $24\cdot 15 + 24\cdot 84 + 24\\ = 24 \cdot (15 + 84 + 1) \\= 2400$; \item $4\cdot 5^2 - 3\cdot 8 + 16:4 + 2017^0 \\= 4\cdot 25 - 24 + 4 + 1 \\= 81$. \end {listEX}  \hfill \qedEX 
\end{Solbook}
\begin{Solbook}{2}
 \par \noindent  \begin {listEX}[3] \item $4\cdot (x + 2) = 24$\\ \hspace *{0.85 cm}$\begin {aligned} x + 2 &= 6\\ x &= 4 \end {aligned}$ \item $-x + 20 = 8$\\ \hspace *{1.25 cm}$\begin {aligned} x &= 20-8\\ x &= 12 \end {aligned}$ \item $x\,\vdots \,5 \Rightarrow x \in B(5)$\\ $x \in \{10;15;20;25\}$ \end {listEX}  \hfill \qedEX 
\end{Solbook}
\begin{Solbook}{3}
 \par \noindent  Tổng số tiền Mai đã mua hết là: $5\cdot 14\,000 + 8\cdot 7\,000 + 2\cdot 35\,000 = 196\,000$ (đồng) \\Số tiền Mai được trả lại là: $200\,000 - 196\,000 = 4000$ (đồng)  \hfill \qedEX 
\end{Solbook}
\begin{Solbook}{4}
 \par \noindent  Gọi $x$ $(x\in \mathbb {N})$ là số học sinh đi tham quan. Ta có $x \in BC(40;45)$ và $600 \leqslant x \leqslant 800)$\\ Phân tích các số ra thừa số nguyên tố $40=2^3\cdot 5 \hspace {1cm}45=3^3\cdot 5$ \\ BCNN$(40;45) =2^3\cdot 3^3\cdot 5= 360$ \\BC$(40;45) =$ B$(360)=\{0; 360;720;1080;... \}$ \\Mà $x \in N, 600 \leqslant x < 800$ \\Nên $x = 720$\\ Vậy có 720 học sinh đi tham quan.  \hfill \qedEX 
\end{Solbook}
\begin{Solbook}{5}
 \par \noindent  \begin {enumerate} \item Diện tích ao nuôi cá nhà bạn Hà: $6\cdot 6 = 36$ (m$^2$) \item Diện tích mảnh vườn nhà bạn Hà là: $25\cdot 18 = 450$ (m$^2$) \\ Diện tích phần đất trồng rau là: $450 - 36 = 414$ (m$^2$) \\ Lượng hạt giống cần để gieo hết phần đất trồng rau là: $414\cdot 3 = 1242$ (gam) \end {enumerate}  \hfill \qedEX 
\end{Solbook}
\begin{Solbook}{6}
 \par \noindent  \begin {enumerate} \item Ngày thứ 7 cửa hàng bán được khối lượng táo nhiều nhất. \item Ngày thứ 3 cửa hàng bán được khối lượng táo ít nhất. \item Khối lượng táo cửa hàng bán được trong tuần đó là: $30\cdot 20 + 2\cdot 10 = 620$ (kg) \end {enumerate}  \hfill \qedEX 
\end{Solbook}
\begin{Solbook}{7}
 \par \noindent  $A = 1+2+2^2+2^3+2^4+...+2^{99}$ \\$A = (1+2+2^2)+2^3(1+2+2^2)+...+2^{96}(1+2+2^2)+2^{97}$ \\ $A = 7 + 2^3\cdot 7+...+2^{96}\cdot 7+2^{97}$ \\ $A = 7(1+2^3+...+2^{96})+2^{97}$ \\Mà $7(1+2^3+...+2^{96})\vdots 7$ và $2^{97}\not \vdots 7$ nên $A\, \not \vdots \, 7$  \hfill \qedEX 
\end{Solbook}
